\documentclass{ctexart}
\usepackage{amsmath, amssymb}
\usepackage{geometry}
\geometry{a4paper, margin=2cm}

\title{\textbf{第9章 静电场} --- 例题}
\author{}
\date{}

\begin{document}
\maketitle

\begin{enumerate}

\item 已知两杆电荷线密度为 $\lambda$，长度为 $L$，相距 $L$。求两带电直杆间的电场力。

\item 在带电量为 $Q$ 的点电荷产生的静电场中，有一带电量为 $q$ 的点电荷。求 $q$ 在 $a$ 点和 $b$ 点的电势能（选无穷远为零势能点）。

\item（例1）两块等面积的金属平板，分别带电荷 $q_A$ 和 $q_B$，平板面积均为 $S$，两板间距离为 $d$，且满足面积的线度远大于 $d$。求静电平衡时两金属板各表面上的电荷面密度。

\item（书中例题8.15，p.314）求均匀电场中任一点的电势及任意两点间的电势差。

\item（例2）半径为 $R_1$ 的导体球带有电荷 $q$，球外有一个内、外半径为 $R_2$、$R_3$ 的同心导体球壳，壳上带有电量 $Q$。求：
\begin{enumerate}
  \item 两球的电势 $V_1$ 和 $V_2$；
  \item 两球的电势差 $\Delta V = V_1 - V_2$；
  \item 用导线把球和球壳联在一起后 $V_1$、$V_2$ 及 $\Delta V$；
  \item 若外球接地，$V_1$、$V_2$ 及 $\Delta V$；
  \item 设外球离地面很远，若内球接地，$V_1$、$V_2$、$\Delta V$ 各为多少？
\end{enumerate}

\item 求电偶极子在延长线上和中垂线上一点产生的电场强度。

\item（书中例题8.18，p.316）均匀带电圆环半径为 $R$，电荷线密度为 $\lambda$。求圆环轴线上一点的电势。

\item 半径为 $R$ 的均匀带电细圆环，带电量为 $q$。求圆环轴线上任一点 $P$ 的电场强度。

\item（书中例题8.19，p.316）计算半径为 $R$、均匀带电量为 $q$ 的圆形平面板轴线上任意一点的电势。

\item（书中例题8.21，p.319）半径为 $R$、带电量为 $q$ 的均匀带电球面。试求球内外任意一点的电势。

\item 面密度为 $\sigma$ 的圆板在轴线上任一点的电场强度。

\item 半径为 $R$、带电量为 $q$ 的均匀带电球体。求带电球体的电势分布。

\item（例3）接地导体球附近有一点电荷，求导体上的感应电荷。

\item（例8.22）"无限长"均匀带电圆柱面的半径为 $R$，单位长度上带电量为 $\lambda$，试求其电势分布。

\item（教材例题8.36，p.346）平行板电容器，其中充有两种均匀电介质。求：
\begin{enumerate}
  \item 各电介质层中的场强；
  \item 极板间电势差和电容。
\end{enumerate}

\item（书中例题8.37，p.347）半径分别为 $R_1$ 和 $R_3$ 的同心导体球面组成的球形电容器，中间充满相对介电常数为 $\varepsilon_{r1}$ 和 $\varepsilon_{r2}$ 的两层各向同性均匀介质，它们的分界线为 $R_2$ 的同心球面。求此电容器的电容。

\item（书中例题8.32，p.335）由两个半径为 $a$ 的平行长直导线，轴间距离为 $d \gg a$，线电荷密度分别为 $+\lambda$ 和 $-\lambda$。求单位长度平行直导线之间的电容。

\item（补充例题）同心导体球面组成的球形电容器，半径分别为 $R_1$ 和 $R_2$，带电量分别为 $\pm Q$，电容器下半部充有电介质油，相对电容率为 $\varepsilon_r$。求：(1) 介质中任意点的电场强度和电位移矢量；(2) 电容器的电容。

\item 长为 $L$ 的均匀带电直杆，电荷线密度为 $\lambda$。求它在空间一点 $P$ 产生的电场强度（$P$ 点到杆的垂直距离为 $a$）。

\item（书中例题8.38，p.348）平行板电容器的极板面积为 $S$，极板间距 $d$，中间充满相对介电常数为 $\varepsilon_r$ 的电介质。当充电后，两极板间的电势差为 $\Delta u$。求：
\begin{enumerate}
  \item 电容器中电场的能量；
  \item 如果切断充电电源，把电介质从电容器中抽出来，外界要作多少功。
\end{enumerate}

\item（书中例题8.33，p.338）半径为 $a$、带电量为 $q$ 的孤立金属球。求它所产生的电场储存的静电能。

\item（书中例题8.39，p.349）球形电容器中充满了相对介电常数为 $\varepsilon_r$ 的各向同性均匀介质。给电容充电，使其两极上带电量为 $\pm q$。求电容器中电场的能量。

\item（书中例题8.35，p.339）如图两电容并联，$C_1 = 1\,\mu\text{F}$，$u_1 = 100\,\text{V}$；$C_2 = 2\,\mu\text{F}$，$u_2 = 200\,\text{V}$。求并联前后电容器所储存的静电能。

\item 已知"无限长"均匀带电直线的电荷线密度为 $+\lambda$。求距直线 $r$ 处一点 $P$ 的电场强度。

\item 均匀带电球面，总电量为 $Q$，半径为 $R$。求电场强度分布。

\item 已知球体半径为 $R$，带电量为 $q$（电荷体密度为 $\rho$）。求均匀带电球体的电场强度分布。

\item（书中例题8.11，p.304）已知"无限大"均匀带电平面上电荷面密度为 $\sigma$。求电场强度分布。

\end{enumerate}

\end{document}
